% from ega4-errata-addenda-list3-fr.tex
\oldpage[IV]{356}
de $\operatorname{Spec}(A)$; puisque $A$ est réduit, il est donc
intègre. Soient $K$ le corps des fractions de $A$, $A'$ la clôture
intégrale de $A$; par platitude, $B\otimes_AA'$ s'identifie à un
sous-anneau de $B\otimes_AK$ contenant $B$, et comme les éléments non
nuls de $A$ sont $B$-réguliers par platitude, $B\otimes_AK$ s'identifie
à un sous-anneau de l'anneau total des fractions $R$ de $B$. Pour
montrer que $A$ est unibranche, il suffit de voir que $A'$ ne peut avoir
deux idéaux maximaux distincts (nécessairement au-dessus de $y$
(Bourbaki, \emph{Alg. comm.}, chap.~V, \S2,
n\textsuperscript{o}~1, prop.~1)); s'il en était ainsi,
$B'=B\otimes_AA'$ aurait deux idéaux premiers distincts au-dessus de
$x$ (I, 3.4.7), et il suffit de voir que ce n'est pas possible.

Soit $\mathfrak N$ le nilradical de $B$; si $S$ est l'ensemble des
éléments réguliers de $B$, on a $R=S^{-1}B$, et
$S^{-1}\mathfrak N=R\mathfrak N$ est le nilradical de $R$, donc
\[
 R/R\mathfrak N=S^{-1}B/S^{-1}\mathfrak N=S^{-1}(B/\mathfrak N)
\]
est contenu dans le corps des fractions $L$ de l'anneau local intègre
$B/\mathfrak N=B_{\mathrm{red}}$; il en est donc de même de
\[
 B'_{\mathrm{red}}=B'/(B'\cap R\mathfrak N)
\]
qui est entier sur $B_{\mathrm{red}}$; par hypothèse
$B'_{\mathrm{red}}$, étant contenu dans la clôture intégrale de
$B_{\mathrm{red}}$, est un anneau local, donc il en est de même de
$B'$, ce qui établit notre assertion.

Si l'on suppose en outre $B$ géométriquement unibranche, le corps
résiduel de $B'_{\mathrm{red}}$, étant contenu dans le corps résiduel de
la clôture intégrale de $B_{\mathrm{red}}$, est une extension radicielle
du corps résiduel $k(x)$ de $B$. Notons maintenant que le corps résiduel
de $B'_{\mathrm{red}}$ est aussi celui de $B'=B\otimes_AA'$, donc est
égal à $k(x)\otimes_{k(y)}k'$, où $k'$ est le corps résiduel de $A'$;
on déduit alors de {\hyperref[canon-index.ega4.c04.prop.2-6-1]{(2.6.1)}} que $k'$
est une extension radicielle de $k(y)$, donc $A$ est aussi
géométriquement unibranche.

\paragraph{(Err\textsubscript{IV}, 33).}
Dans ({\hyperref[canon-index.ega4.c04.item.12-1-2]{IV, 12.1.2}}), l'exemple donné
en (i) n'est pas correct, le morphisme considéré étant ouvert. On
construit un exemple d'un morphisme plat $f:X\to Y$ dont la restriction
à une composante irréductible de $X$ n'est pas un morphisme ouvert de
la façon suivante. Soit $Z=\operatorname{Spec}(\mathbb C[U,V])$ le
plan affine sur le corps des nombres complexes, et considérons le
schéma $Y$ obtenu en recollant deux points fermés distincts $a$, $b$ de
$Z$; on a $Y=\operatorname{Spec}(B)$, où $B$ est le sous-anneau de
$\mathbb C[U,V]$ formé des polynômes prenant la même valeur aux points
$a$ et $b$ (l'ensemble des points fermés de $Z$ étant identifié à
$\mathbb C^2$); il est immédiat que $B$ est un sous-$\mathbb C$-espace
vectoriel de $\mathbb C[U,V]$ de codimension 1, donc $\mathbb C[U,V]$
est un $B$-module de type fini, et par suite le morphisme $Z\to Y$ est
fini, l'image réciproque du point $c$ de $Y$ obtenu par recollement de
$a$ et de $b$ étant l'ensemble $\{a,b\}$.

Considérons alors le schéma $X$ obtenu en recollant deux exemplaires de
$Z$, soient $Z_1$, $Z_2$, de sorte que si $a_i$, $b_i$ ($i=1,2$) sont
les points de $Z_i$ correspondant à $a$ et $b$, $a_1$ soit recollé à
$b_2$ et $b_1$ à $a_2$ (cf. ({\hyperref[canon-index.ega4.c04.item.6-5-5]{6.5.5}},
(ii))). On voit aisément (\emph{loc. cit.}) que le morphisme $X\to Y$
ainsi obtenu est plat et fini (et même étale) et que $Z_1$ et $Z_2$
s'identifient aux deux composantes irréductibles de $X$. Cependant, si
$D$ est une droite dans $Z_1$ passant par $a_1$ mais non par $b_1$,
l'image $V$ dans $Y$ de l'ouvert $Z_1-D$ de $Z_1$ contient l'image de
$b_1$, c'est-à-dire le point $c$, mais ne contient pas l'image du point
générique de $D$, donc $V$ n'est pas un voisinage de $c$.

\paragraph{(Err\textsubscript{IV}, 34).}
Dans le titre de (IV, 13.2), ligne 5 du bas de la p.~190, remplacer :
« équidimensionnels » par « localement équidimensionnels ». De même
dans (IV, 13.2.2), ligne 4 de la p.~191, remplacer deux fois
« équidimensionnel » par « localement équidimensionnel ».

